\documentclass[11pt,a4paper]{article}

\usepackage[T1]{fontenc}
\usepackage[utf8]{inputenc}
\usepackage{lmodern}
\usepackage{microtype}
\usepackage[margin=2.5cm]{geometry}
\usepackage{amsmath,amssymb}
\usepackage{booktabs,tabularx,multirow}
\usepackage[group-separator={,},group-minimum-digits=4]{siunitx}
\usepackage{xcolor}
\usepackage[font=small,labelfont=bf]{caption}
\usepackage[numbers,sort&compress]{natbib}
\usepackage[colorlinks=true,linkcolor=teal!70!black,citecolor=teal!70!black,urlcolor=teal!70!black]{hyperref}
\usepackage{cleveref}

\newcommand{\ioVol}{\iota_j}
\newcommand{\tauIO}{\tau^{\mathrm{IO}}}
\newcommand{\kapIO}{\kappa^{\mathrm{IO}}}
\newcommand{\kapVM}{\kappa^{\mathrm{VM}}}
% typewriter text that may break after an underscore
\newcommand{\codeus}{\textunderscore\allowbreak}
\newcommand{\code}[1]{{\ttfamily\let\_\codeus #1}}
\setlength{\emergencystretch}{3em}

\title{Experimental setup:\\the MOSPACE generated instances}
\author{MOSPACE project --- methodology note}
\date{20 November 2026}

\begin{document}
\maketitle

\begin{abstract}
This note describes the generated benchmark used to evaluate the MOSPACE
cloud-bursting scheduler: how its instances are built, what each experimental
factor controls, and the protocol under which exact and heuristic methods are
compared on it. The benchmark is a full factorial of 108 instances: 2 machine
catalogues $\times$ 2 real Slurm workloads $\times$ 3 on-prem node counts
$\times$ 3 initial cluster loads $\times$ 3 result-data levels. Jobs come
unchanged from two snapshots of a production reservoir-simulation cluster (964 and
\num{3340} jobs). Cloud types are priced from AWS and Azure on-demand list prices,
and output volumes come from the model in the companion note on result data. The
instance name determines the instance, so the generator rebuilds every instance
bit for bit; all 108 files match their recorded sha256.
\end{abstract}

\section{Scheduling problem}\label{sec:problem}
Each instance is an offline bi-objective scheduling problem on unrelated parallel
machine \emph{types}. A finite number of on-prem nodes is available, together with
an unlimited, on-demand cloud. A job $j$ has a release slot $a_j$, an execution
time $e_j$ (slots, fractional), a core count $c_j$, a memory requirement and an
output volume $\ioVol$ (GB). It runs on one node of a type $m$ that covers its
cores and memory, and all jobs have unit weight. For job $j$ on type $m$, with
relative performance $P_m$ (equal to 1 for every type here), the billed time, the
occupied time and the cost are
\begin{align}
  p^{\mathrm{bill}}_{jm} &= e_j / P_m + \tauIO_m\,\ioVol, &
  p^{\mathrm{occ}}_{jm}  &= \max\bigl(1, \lceil p^{\mathrm{bill}}_{jm}\rceil\bigr), &
  c_{jm} &= \kapIO_m\,\ioVol + \kapVM_m\,p^{\mathrm{bill}}_{jm}.
\end{align}
The two objectives are
\begin{equation}
  f_1 = \sum_j (C_j - a_j) \quad\text{(total turnaround, slots)}, \qquad
  f_2 = \sum_j c_{j m(j)} \quad\text{(total cost, USD)},
\end{equation}
where $C_j$ is the completion slot. On-prem types have
$\kapVM_m=\kapIO_m=\tauIO_m=0$, so on-prem work costs nothing, but on-prem
capacity is limited and partly held by jobs already running at $t=0$. Cloud
work costs money and, through $\tauIO_m$, holds its VM while the results are
uploaded. The slot length is \SI{600}{s} and the horizon is $H=289$
(W1) or 288 (W2) slots, i.e.\ about two days.

\section{Workloads}\label{sec:workloads}
The jobs, release times, runtimes, resources, horizon and slot length are copied
unchanged from two real snapshots (\cref{tab:workloads}). Runtimes are the
observed elapsed times (the known-runtime scenario). The original cloud types
and the original running jobs of each snapshot are replaced by the generated ones
(\cref{sec:catalogue,sec:factors}). The end times of the original running jobs
serve only as the empirical distribution from which the generator draws.

\begin{table}[ht]
\centering\small
\caption{The two source workloads. Execution times are in 600-s slots;
``ensemble'' counts jobs that share their (user, job-name stem) with at least
10 jobs (\cref{sec:io}).}
\label{tab:workloads}
\begin{tabularx}{\linewidth}{lXX}
\toprule
 & W1 & W2 \\
\midrule
Snapshot & window of 2024-11-25, 21:17 UTC & congestion event of 2024-01-30, 20:00--24:00 UTC \\
Jobs & 964 & \num{3340} \\
Release window & slots 0--10.8 & slots 0--24.0 \\
Cores (min / median / max) & 20 / 40 / 64 & 1 / 30 / 64 \\
Memory, GB (min / median / max) & 150 / 300 / 480 & 7.5 / 225 / 480 \\
Execution time, slots (min / median / max) & 0.05 / 3.7 / 68.2 & 0.003 / 5.2 / 268.4 \\
Total work, core-hours & \num{44177} & \num{215718} \\
Ensemble members & 945 (98\,\%) & \num{3194} (96\,\%) \\
Horizon $H$ & 289 slots & 288 slots \\
Running-job end-time pool & \num{1654} values, median 21 & \num{1525} values, median 23 \\
\bottomrule
\end{tabularx}
\end{table}

\section{Machine catalogue}\label{sec:catalogue}
Two catalogues (\emph{classes}) define the on-prem node types and their cloud
equivalents (\cref{tab:catalogue}). Class~A models the two real node types of the
cluster (csr1: 40 cores / 300 GB, csr3: 64 cores / 480 GB), each with one AWS
equivalent. Class~B has four on-prem types, each with an AWS and an Azure
equivalent, which lets the scheduler trade one provider against the other.

A cloud equivalent is the smallest instance of the same Intel generation (Ice
Lake: AWS m6i, Azure Dsv5) whose physical cores (vCPU/2, since SMT is on) and
memory cover the on-prem node. Azure has no 128-vCPU Dsv5, so the 64-core node maps
to D128s\_v6. Prices are Linux on-demand list prices (AWS us-east-1, Azure eastus,
\code{data/cloud\_prices.json}), converted to a cost per slot
$\kapVM_m = \text{USD/h}\times 600/3600$. All types run at $P_m=1$ with no
deployment delay, and every cloud type has unlimited capacity.

\begin{table}[ht]
\centering\small
\caption{On-prem node types and their cloud equivalents (USD per hour).}
\label{tab:catalogue}
\begin{tabular}{lllcc}
\toprule
On-prem type & AWS & Azure & Class A & Class B \\
\midrule
c16 (16 c / 128 GB) & m6i.8xlarge (1.536)  & D32s\_v5 (1.536)  & -- & W2 only \\
c32 (32 c / 256 GB) & m6i.16xlarge (3.072) & D64s\_v5 (3.072)  & -- & \checkmark \\
c40 (40 c / 300 GB) & m6i.24xlarge (4.608) & D96s\_v5 (4.608)  & \checkmark\ (AWS only) & \checkmark \\
c64 (64 c / 480 GB) & m6i.32xlarge (6.144) & D128s\_v6 (6.451) & \checkmark\ (AWS only) & \checkmark \\
\bottomrule
\end{tabular}
\end{table}

\paragraph{Pickable types only.} A type that no job of the workload would ever
use is dropped from that workload's instances. This covers an on-prem type that no
job fits (no W1 job fits c16, since its smallest job has 20 cores) and
a cloud size that fits no job or is strictly dearer, for every job it fits, than
another size of the \emph{same} provider. Dominance is checked within each
provider, so each provider keeps a complete offer: D128s\_v6 stays even though it
is dearer than m6i.32xlarge. Class~A therefore always has 2 on-prem and 2 cloud
types. Class~B has 3 on-prem and 6 cloud types for W1, and 4 and 8 for
W2. A memory-heavy 32c/512\,GB type was considered for class~B and
removed: it fit exactly the same jobs as c32 and only split c32's capacity.

\paragraph{Data transfer.} Cloud jobs pay internet egress
$\kapIO_m$ = \$0.09/GB (AWS) or \$0.08/GB (Azure, routing preference: internet),
both at the first paid tier, North America. The transfer time follows the
minimum-of-subsystems bound of \citet{liu2017explaining}: an upload is no faster
than its slowest stage among source-disk read, network and sink write. The
provider network caps are at least \SI{5}{Gbps} on every VM used here
\citep{aws-bandwidth,azure-bandwidth}, an AWS egress cap that
\citet{jain2023skyplane} also measured. The binding stage is therefore the
throughput each provider includes in the price of its block storage: 125\,MiB/s
($\approx$\SI{131}{MB/s}) for an EBS gp3 volume \citep{aws-ebs} and \SI{125}{MB/s}
for an Azure Premium SSD v2 disk \citep{azure-disks}. On Azure this assumes the
results are written to a Premium SSD v2 data disk, which cannot serve as the OS
disk. These baselines give $\tauIO_m$ = \SI{7.63}{s/GB} on AWS (0.01272 slot/GB)
and \SI{8.0}{s/GB} on Azure (0.01333 slot/GB).

The two rates are nearly equal. Their 5\,\% difference is not a measured
difference in transfer speed: both providers quote a baseline of 125, but AWS
quotes it in MiB/s and Azure in MB/s. Neither rate was measured end to end; both
are documented values. In the instances, the difference makes Azure the cheaper
but marginally slower egress, so neither provider dominates the other on both
objectives. This trade-off should be read as an assumption of the model rather
than as an observed property of the providers. The derivation is in
\code{docs/result\_data\_model.tex}.

\section{Experimental factors}\label{sec:factors}
Every factor takes one level for the whole instance (\cref{tab:factors}). The
full factorial has $2\times3\times3\times3=54$ instances per class and 108 in all.
An instance is named by its factors, e.g.\ \code{B\_W2\_n300\_lE\_iM} for class~B,
the W2 workload, $N=300$, load E and io level M.

\begin{table}[ht]
\centering\small
\caption{Factors and levels.}
\label{tab:factors}
\begin{tabularx}{\linewidth}{llX}
\toprule
Factor & Levels & Meaning \\
\midrule
Class & A, B & machine catalogue (\cref{sec:catalogue}) \\
Workload & W1, W2 & job set (\cref{sec:workloads}) \\
Node count $N$ & 100, 300, 500 & nodes of \emph{every} on-prem type (any integer in $[100,500]$ is
  supported; the factorial uses these three) \\
Load & F, H, E & jobs running at $t=0$ on each on-prem type: 0 (free),
  $\lfloor N/2\rfloor$ (half) or $N$ (entirely full) \\
io level & L, M, H & output-retention policy that sets $\ioVol$ (\cref{sec:io}) \\
\bottomrule
\end{tabularx}
\end{table}

\paragraph{Initial load.} Each running job holds one node of its type until its
end slot. End slots are drawn with replacement from the workload's own empirical
distribution of running-job end times (\cref{tab:workloads}). In W2, a
quarter of these ends fall within the first hour, while about 2\,\% lie beyond the
horizon, so those nodes stay blocked for the whole instance. The random number
generator is seeded by sha256(class, workload, type, $N$) and not by the load.
The $N$ draws are made once, and load H takes the first $\lfloor N/2\rfloor$ of
them, so the half-full set is a prefix of the full one and raising the load only
ever removes capacity.

\paragraph{Why one level per instance.} An earlier design set the count (and
load) separately for each on-prem type. The space then grows as $3^k$ in the number
of types $k$ (\num{1944} class~B instances, \num{36450} with per-type loads),
yet the jobs can hardly tell those instances apart: the fit sets are nested
(c16 $\subset$ c32 $\subset$ c40 $\subset$ c64), and only 5 W1 and 26
W2 jobs need c64 rather than c40. A single count, load and io level per
instance keeps the space small enough to run both classes as full factorials.

\section{Result-data volume}\label{sec:io}
The Slurm logs record no I/O, so $\ioVol$ is modelled (full derivation and
evidence in \code{docs/result\_data\_model.tex}):
\begin{equation}
  \ioVol = r \cdot 40 \cdot (c_j/40)^{0.54} \cdot h_j \cdot 10^{0.5\,z_j},
  \qquad z_j\sim\mathcal N(0,1)\text{ clipped to }\pm3,
\end{equation}
with $h_j$ the runtime in hours and $r$ the job's output-retention rate in GB per
core-hour: S = 0.02 (summary vectors only), R = 0.13 (grid, init and a final
restart) or F = 1.5 (a restart at every report step). These rates were measured
with OPM Flow and Darshan on open reservoir decks. Each job is classified as an
\emph{ensemble member} (history-matching or optimisation runs) or a \emph{single
run}, and the io level assigns one policy to each class:
\begin{center}\small
\begin{tabular}{llll}
\toprule
Level & Ensemble & Single & Reading \\
\midrule
L & S & R & results stay in the cloud, only curves return \\
M & R & F & production decks as shipped \\
H & F & F & full restart history brought home \\
\bottomrule
\end{tabular}
\end{center}
Each $z_j$ is seeded per (workload, job), so every level uses the same draw (common
random numbers) and $\ioVol$ only grows from L to H. A horizon guard caps
$\ioVol$ at $(H-\lceil a_j\rceil-e_j)/\tau_{\max}$, with $\tau_{\max}$ the slowest
cloud $\tauIO$ of the instance, so that every cloud type stays feasible for a job
started at its release. The guard binds only on W2: on 3 of \num{3340}
jobs at level M and 7 at level H. The resulting total volumes are in
\cref{tab:instances}.

\section{Resulting instances}\label{sec:instances}
\Cref{tab:instances} summarises the 108 instances by cell. Within a
(class, workload) pair, the jobs and types are fixed. The count and load set the
on-prem capacity and its occupancy, and the io level sets the volumes.

\begin{table}[ht]
\centering\small
\caption{Instance dimensions. On-prem nodes $= N\times$ (number of on-prem types);
running jobs at $t=0$ $=$ load fraction $\times$ on-prem nodes. Total
$\sum_j\ioVol$ in TB per io level L / M / H (the W2 M and H totals vary by
under 0.5\,\% with the class, through $\tau_{\max}$ in the guard).}
\label{tab:instances}
\setlength{\tabcolsep}{4pt}
\begin{tabular}{llrcccc}
\toprule
Class & Workload & Jobs & Types (on-prem / cloud) & On-prem nodes & Instances & $\sum\ioVol$ (TB) \\
\midrule
A & W1 & 964 & 2 / 2 & 200 / 600 / \num{1000} & 27 & 2.6 / 20.9 / 142.2 \\
A & W2    & \num{3340} & 2 / 2 & 200 / 600 / \num{1000} & 27 & 16.3 / 147 / 640 \\
B & W1 & 964 & 3 / 6 & 300 / 900 / \num{1500} & 27 & 2.6 / 20.9 / 142.2 \\
B & W2    & \num{3340} & 4 / 8 & 400 / \num{1200} / \num{2000} & 27 & 16.3 / 147 / 639 \\
\bottomrule
\end{tabular}
\end{table}

\paragraph{Trivial instances.} When there is enough free on-prem capacity, every
job can start on-prem at its release, so the fastest schedule is also free and the
Pareto front is a single point ($f_2^T=0$; see \cref{sec:protocol}). This happens
in 21 of the 108 instances: all W1 instances with $N=300$ at load F or
$N=500$ at load F or H, in both classes (18), and \code{B\_W2\_n500\_lF} at all
three io levels (3). They are reported separately and excluded from quality
statistics.

\paragraph{Monotonicity.} Since the running-job sets are nested across loads and
volumes across io levels, the optimal $f_1$ and $f_2$ can only grow along a load
chain F$\to$H$\to$E or an io chain L$\to$M$\to$H. The test suite checks these
properties with exact solves.

\paragraph{Reproducibility.} The generator (modules \code{instance\_gen} and
\code{io\_model}) is deterministic: its inputs are the instance name,
\code{data/cloud\_prices.json}, the two workload files and
\code{data/job\_classes\_<workload>.json}. The instance files are written by
\begin{center}
\code{uv run python scripts/build\_generated\_instances.py}
\end{center}
into \code{data/generated/}, and \code{manifest.json} there records each file's
sha256. All 108 files were rebuilt on 2026-11-03 and again on 2026-11-20, and every
hash matched.

\section{Evaluation protocol}\label{sec:protocol}
Per instance, the reference benchmark (\code{scripts/\allowbreak bench\_generated.py}, results
in \code{results/\allowbreak generated\_bench/\allowbreak bench\_io\_noreduc.jsonl}) runs the following
methods in order.

\begin{enumerate}
\item \textbf{Lexicographic anchors (Gurobi).} lex-$(f_1,f_2)$ minimises $f_1$, then
  $f_2$ with $f_1$ fixed, giving the ideal-turnaround point $(f_1^T, f_2^T)$.
  lex-$(f_2,f_1)$ minimises $f_2$, then $f_1$, giving $(f_1^0, f_2^0)$; here
  $f_2^0=0$ in every instance, because the on-prem-only schedule is feasible.
\item \textbf{Baselines.} Greedy earliest completion (completion first, then
  cost, least slack first); on-prem only (SPT, never bursts; it completes all
  jobs in all 108 instances); and threshold bursting, where a job runs on-prem if
  its earliest on-prem wait is at most $\theta$ slots and otherwise on the cheapest
  cloud type, for $\theta\in\{0,1,2,3,6,12,24,48,96\}$.
\item \textbf{Weighted sums,} for $\lambda\in\{0.25,0.5,0.75\}$, of the
  normalised scalar
  \begin{equation}
    g_\lambda = \lambda\,\frac{f_1-f_1^T}{f_1^0-f_1^T} + (1-\lambda)\,\frac{f_2}{f_2^T},
    \qquad\text{subject to } f_1\le f_1^0,\ f_2\le f_2^T,
  \end{equation}
  solved exactly by Gurobi and heuristically by the dispatch metaheuristic
  (\code{weighted\_solve\_dispatch}). For threshold bursting, the reported $\theta$
  is the one with the lowest $g_\lambda$ at that $\lambda$, chosen after the fact,
  which is an oracle choice in the baseline's favour. On a degenerate instance
  ($f_2^T=0$), the scalar reduces to minimising $f_2$ subject to $f_1\le f_1^0$.
\end{enumerate}
Every schedule is verified with \code{Instance.check\_schedule}. The quality of a
method is its gap $100\,(g_\lambda^{\text{method}}-g_\lambda^*)/g_\lambda^*$ to the
Gurobi value. Because the three io levels of a cell share their jobs, capacity and
$z_j$, io effects are compared as paired ratios within a cell.

\paragraph{Formulation and solver settings.} Gurobi solves the space-time network
formulation (\code{SpaceTimeFormulation}) without its dominance reductions, so the
times are an upper bound on the production configuration. Parameters:
\code{Threads}=4, \code{TimeLimit}=\SI{600}{s} per \code{optimize()} call,
\code{MIPGap}=0. The lex phases also set \code{MIPGapAbs} to 0.999 of one objective
step (1 slot for $f_1$, $\delta=10^{-4}$ USD for $f_2$). Phase~2 of lex-$f_1$ runs
with \code{Presolve}=0, because single-threaded presolve on the dense $f_1$-fixing
row otherwise stalls past the time limit. A lex result counts as optimal only when
both phases are certified. Wall time is measured with \code{time.perf\_counter()}
around each call and includes model build. In this setting, 530 of 540 Gurobi
solves are proven optimal. The 10 exceptions are 8 lex-$f_2$ solves on
W2 instances with $N=100$ at loads H/E, and 2 weighted sums at
$\lambda=0.25$ (\code{B\_W2\_n100\_lE\_iM} and \code{A\_W2\_n100\_lE\_iL}).

\paragraph{Hardware and software.} Intel Core i9-13900 (8 P- and 16 E-cores, 32
threads, up to 5.6\,GHz), 62\,GB RAM, CachyOS Linux (kernel 7.3), CPython 3.14,
gurobipy 13.0.1 (academic licence), numpy 2.4.4. Three instances run in parallel
with 4 threads each, and threads are not pinned to cores. Timings are therefore
compared between methods only through medians.

\section{Limitations}\label{sec:limits}
The jobs are real but come from two snapshots of one cluster, and the factorial
varies the platform around them rather than the jobs. Output volumes are modelled
rather than observed, and the io level expresses that uncertainty. All types run
at $P_m=1$, so on-prem and cloud differ only in cost, capacity and transfer time,
not in speed. Cloud capacity is unlimited, with no quotas, spot pricing or volume
discounts. Uploads are assumed not to contend on the site's WAN link.

\begin{thebibliography}{9}\small
\bibitem[Liu et~al.(2017)]{liu2017explaining} Z.~Liu, P.~Balaprakash, R.~Kettimuthu, and I.~Foster. Explaining wide area data transfer performance. In \emph{Proc. HPDC}, pages 167--178, 2017.
\bibitem[Jain et~al.(2023)]{jain2023skyplane} P.~Jain, S.~Kumar, S.~Wooders, S.~G. Patil, J.~E. Gonzalez, and I.~Stoica. Skyplane: optimizing transfer cost and throughput using cloud-aware overlays. In \emph{Proc. NSDI}, 2023.
\bibitem[AWS(2026a)]{aws-ebs} Amazon Web Services. Amazon EBS General Purpose SSD volumes. \url{https://docs.aws.amazon.com/ebs/latest/userguide/general-purpose.html}.
\bibitem[AWS(2026b)]{aws-bandwidth} Amazon Web Services. Amazon EC2 instance network bandwidth. \url{https://docs.aws.amazon.com/AWSEC2/latest/UserGuide/ec2-instance-network-bandwidth.html}.
\bibitem[Microsoft(2026a)]{azure-disks} Microsoft. Azure managed disk types. \url{https://learn.microsoft.com/en-us/azure/virtual-machines/disks-types}.
\bibitem[Microsoft(2026b)]{azure-bandwidth} Microsoft. Virtual machine network bandwidth. \url{https://learn.microsoft.com/en-us/azure/virtual-network/virtual-machine-network-throughput}.
\end{thebibliography}

\end{document}
